%%%%%%%%%%  scan idx 14  |  folio 1  %%%%%%%%%%
\SGAPrintPage{1}
\sgaStableTarget{sga6.e00.expose}\section*{\underline{Exposé 0}}
\begin{center}
\underline{\underline{OUTLINE OF A PROGRAMME FOR A THEORY OF INTERSECTIONS}}\\[3pt]
\underline{\underline{ON GENERAL SCHEMES}}\\[8pt]
\underline{by A. Grothendieck}
\end{center}

The present exposé is introductory in nature, and reading it is not logically indispensable for the study of the Seminar. It is addressed more particularly to readers familiar with the provisional version of the Riemann--Roch theorem, as it is presented in the report of Borel--Serre [2] or in the report of Grothendieck already mentioned in the Introduction (subsequently cited as [RRR]), which is reproduced as an appendix at the end of the present exposé.

\sgaStableTarget{sga6.e00.sec.1}\subsection*{1.}
Recall the Riemann--Roch formula for a proper morphism
\[
f:X\longrightarrow Y
\]
of quasi-projective smooth schemes over a field $k$, and a coherent sheaf $F$ on $X$, whose class in the group $K(X)$ of classes of coherent sheaves on $X$ is denoted by $\cl(F)$:
\sgaStableTarget{sga6.e00.eq.1-1}\begin{equation}\tag{1.1}
\Todd(T_Y)\,\ch_Y\bigl(f_*(\cl(F))\bigr)=f_*\bigl(\Todd(T_X)\,\ch_X(F)\bigr).
\end{equation}
This formula holds in $A(Y)\otimes_{\mathbb Z}\mathbb Q$, where $A(Y)$ is the Chow ring of $Y$, the $f_*$ on the right-hand side being obtained by tensoring with $\mathbb Q$ the homomorphism ``direct image of cycles''
\[
f_*:A(X)\longrightarrow A(Y),
\]
and that on the left-hand side being the Euler--Poincaré characteristic of $F$ relative to $f$:
\[
f_*(\cl(F))=\sum_i(-1)^i\cl\bigl(R^if_*(F)\bigr)~;
\]
$\ch_X$, $\ch_Y$ denote the Chern characters on $X$, respectively on $Y$, and $T_X$, $T_Y$ the tangent bundles of $X$, respectively of $Y$. As is known, $\Todd(-)$ and $\ch(-)$ are universal polynomials with coefficients in $\mathbb Q$ in the Chern classes of the argument. Since the constant term of $\Todd(-)$ is $1$, it is an invertible element for every value of the argument, so that, after multiplication by $\Todd(T_Y)^{-1}$, formula (1.1) takes the equivalent form, more convenient for our purpose,
\sgaStableTarget{sga6.e00.eq.1-2}\begin{equation}\tag{1.2}
\ch_Y\bigl(f_*(\cl(F))\bigr)=f_*\bigl(\Todd(T_f)\cdot\ch_X(F)\bigr),
\end{equation}
where we have put
\sgaStableTarget{sga6.e00.eq.1-3}\begin{equation}\tag{1.3}
T_f=T_X-f^*(T_Y)\in K(X),
\end{equation}

%%%%%%%%%%  scan idx 15  |  folio 2  %%%%%%%%%%
\SGAPrintPage{2}
so that $T_f$ plays the role of a \underline{virtual relative tangent bundle} of $X$ over $Y$. When the morphism $f$ is smooth (i.e. its tangent map is everywhere surjective), we simply have
\[
T_f=T_{X/Y}\qquad\text{(tangent bundle along the fibres)},
\]
whereas when $f:X\to Y$ is an immersion, we find
\[
T_f=-\check N_{X/Y},
\]
where $\check N_{X/Y}$ denotes the normal sheaf of $X$ in $Y$.

One of the main aims of the Seminar is to generalize (1.2) simultaneously in two directions:
\begin{enumerate}[label=\alph*)]
\item To remove the hypothesis of the existence of a base field $k$.
\item To replace the regularity hypotheses on $Y,X$ by a hypothesis of ``local regularity'' of $f$.
\end{enumerate}
Finally, along the way, we shall also consider the problem:
\begin{enumerate}[label=\alph*),start=3]
\item To eliminate the quasi-projectivity hypotheses which, in the absence of a base field, are expressed by the existence of ample invertible Modules on $X$ and on $Y$.
\end{enumerate}

\sgaStableTarget{sga6.e00.sec.2}\subsection*{2.}
Let us first examine generalization a), while retaining the hypotheses of regularity and of the existence of ample invertible Modules on $X$ and on $Y$.

The definition of $K(X)$, $K(Y)$ and of the homomorphism $f_*:K(X)\to K(Y)$ then presents no new problem, because $X$ and $Y$ are regular. The most natural way of giving a meaning to (1.2) therefore seems to consist in defining \underline{Chow rings} $A(X),A(Y)$ and a group homomorphism
\[
f_*:A(X)\longrightarrow A(Y),
\]
in establishing a theory of Chern classes, providing maps
\[
c_i:K(X)\longrightarrow A(X),
\]
and similarly for $Y$, and finally in describing a virtual relative tangent bundle element
\[
T_f\in K(X).
\]

%%%%%%%%%%  scan idx 16  |  folio 3  %%%%%%%%%%
\SGAPrintPage{3}
\sgaStableTarget{sga6.e00.sec.2.1}\subsubsection*{2.1}
As for the latter, a definition presents itself fairly evidently. We use the fact that, by the hypothesis of the existence of an ample invertible Module on $X$ and on $Y$, the morphism $f$ can be factored as
\sgaStableTarget{sga6.e00.eq.2-1}\begin{equation}\tag{2.1}
X\xrightarrow{i}X'\xrightarrow{f'}Y,
\end{equation}
where $i$ is a closed immersion and $f'$ is a smooth morphism; for example, we may take $X'=\mathbb P^r_Y$ for suitable $r$. We then put
\sgaStableTarget{sga6.e00.opening.eq.2-2}\begin{equation}\tag{2.2}
\check T_f=\cl({\underline\Omega}^1_{X'/Y})-\cl({\underline N}_{X/X'}),
\end{equation}
where ${\underline\Omega}^1_{X'/Y}$ is the locally free Module of relative $1$-differentials of $X'$ over $Y$ (or relative cotangent Module of $X'$ over $Y$), and ${\underline N}_{X/X'}$ is the conormal Module of $X$ in $X'$, equal by definition to ${\underline J}/{\underline J}^2$, where ${\underline J}$ is the Ideal on $X'$ defining the closed subscheme $X$; since $X$ and $X'$ are regular, we also conclude that ${\underline N}_{X/X'}$ is locally free. Finally, $\check T_f$ denotes the dual of $T_f$. We readily verify (Exp. VIII) that the element $T_f$ defined by (2.1) does not depend on the chosen factorization of $f$.

\sgaStableTarget{sga6.e00.sec.2.2}\subsubsection*{2.2.}
As for a definition, under the conditions considered, of a Chow ring and a corresponding theory of Chern classes, although it is plausible that such a theory should be developable on the model of the theory already known for smooth quasi-projective algebraic schemes, it seems that this work has not been done at present, and a fortiori it has not been undertaken in the Seminar. Instead, we introduce another ring which will play a role analogous to that of the Chow ring, and whose tensor product with $Q$ is indeed isomorphic to $A(X)\otimes_{\mathbb Z}\mathbb Q$ when $X$ is quasi-projective and smooth over a field $k$. This ring is the associated graded ring of $K(X)$ for a suitable filtration of that ring. A first natural filtration of $K(X)$, the one considered in [RRR], is the ``\underline{topological filtration}'' by the codimension of the support of coherent sheaves, $\operatorname{Filt}^{i}_{\mathrm{top}}K(X)$ being generated by the $\cl(F)$ with $F$ coherent and $\operatorname{codim}(\operatorname{supp}F,X)\ge i$. A first question is the compatibility of this filtration with the ring structure, a question resolved in [RRR], when $X$ is smooth over a field $k$, by means of Chow's ``moving lemma''. Granting this point, we therefore obtain an associated graded ring, which we shall denote here by $\Gr^{\bullet}_{\mathrm{top}}(X)$. In [RRR] a surjective ring homomorphism is defined:
\[
\phi:A(X)\longrightarrow \Gr^{\bullet}_{\mathrm{top}}(X),
\]

%%%%%%%%%%  scan idx 17  |  folio 4  %%%%%%%%%%
\SGAPrintPage{4}
(sending the class of a prime cycle $Z$ to $\cl(\underline O_Z)$,) whose kernel subsequently turns out (Exp. XIV 4.2) to be a torsion subgroup. Moreover, [RRR] shows how the images in $\Gr^{\bullet}_{\mathrm{top}}(X)$ of the Chern classes $c_i(F)$ can be calculated directly in terms of $\cl(F)\in K(X)$ and the exterior-power operations $\lambda^i$ in $K(X)$.

This shows that, for $X$ a noetherian regular scheme equipped with an ample invertible Module, the ring $\Gr^{\bullet}_{\mathrm{top}}(X)$ is a convenient substitute for the hypothetical Chow ring $A(X)$; it also has the advantage over the latter that certain calculations in [RRR] take place directly in this ring and lead to important formulas in it, formulas which, even for $X$ projective and smooth over a field, have not at present been established in $A(X)$ itself (Exp. XIV, 4.3 and 4.4). To justify this point of view completely, however, one would have to resolve the problem mentioned above of the compatibility of the topological filtration of $K(X)$ with its ring structure, a problem which seems essentially connected with Chow's ``moving lemma'', which is itself the essential technical ingredient of the theory of the Chow ring that we have claimed to be able to avoid.

\sgaStableTarget{sga6.e00.sec.2.3}\subsubsection*{2.3.}
To avoid this problem, we give $K(X)$ a second filtration, finer than the topological filtration, which can be described in purely algebraic terms from the augmented $\lambda$-ring structure of $K(X)$. The definition of this filtration is suggested by the expression given in [RRR] for the Chern classes of an $x\in K(X)$ in $\Gr^{\bullet}_{\mathrm{top}}(X)$,
\sgaStableTarget{sga6.e00.eq.2-3}\begin{equation}\tag{2.3}
c_i(x)=\gamma^i(x-\epsilon(x))\mod \operatorname{Filt}^{i+1}_{\mathrm{top}}K(X),
\end{equation}
where $\gamma^i$ denotes a certain linear combination of the $\lambda^j$ ($j\leq i$) made explicit in loc. cit., $\epsilon$ being the augmentation. We show that $\gamma^i(x-\epsilon(x))$ does have topological filtration $\ge i$, which gives a meaning to formula (2.3). We then define the $\lambda$-\underline{filtration} of $K(X)$ by the ideals $\operatorname{Filt}^i K(X)$, where $\operatorname{Filt}^i K(X)$ consists of sums of expressions of the form
\[
\gamma^{i_1}(x_1-\epsilon(x_1))\cdots \gamma^{i_k}(x_k-\epsilon(x_k)),\qquad i_1+\cdots+i_k\ge i.
\]
The associated graded ring of this filtration will be denoted by $\Gr^\bullet(X)$, and it is this ring that will replace the Chow ring of $X$ in the Seminar. We can moreover show that, in the case considered here, the canonical homomorphism
\sgaStableTarget{sga6.e00.eq.2-4}\begin{equation}\tag{2.4}
\Gr^\bullet(X)\longrightarrow \Gr^{\bullet}_{\mathrm{top}}(X)
\end{equation}

%%%%%%%%%%  scan idx 18  |  folio 5  %%%%%%%%%%
\SGAPrintPage{5}
is an isomorphism modulo torsion, which justifies the point of view that this ring tensored with $\mathbb Q$ can play in every respect the same role as the Chow ring tensored with $\mathbb Q$. On the other hand, we have done exactly what is needed for formula (2.3) (with $\operatorname{Filt}^i_{\mathrm{top}}$ replaced by $\operatorname{Filt}^i$) to define a theory of Chern classes with values in $\Gr^\bullet(X)$, from which the Chern classes in $\Gr^{\bullet}_{\mathrm{top}}(X)$ are deduced using the homomorphism (2.4). We show (Exp. V), using the fact (Exp. VI) that $K(X)$ is a \underline{special} $\lambda$-ring in the sense of [RRR] (i.e. a $\lambda$-ring in the sense of Exp. V), that this notion of Chern class satisfies all the usual formal properties reviewed in [RRR].

\sgaStableTarget{sga6.e00.sec.2.4}\subsubsection*{2.4.}
We note a very appreciable advantage of the definition of $\Gr^\bullet(X)$ over that of $A(X)$ and $\Gr^{\bullet}_{\mathrm{top}}(X)$: it remains meaningful for an absolutely arbitrary scheme $X$, and more generally for an arbitrary locally ringed topos $X$, taking for $K(X)$ the group of classes constructed from locally free Modules on $X$, which is still an augmented $\lambda$-ring, hence equipped with a corresponding filtration allowing an associated graded $\Gr^\bullet(X)$ to be defined. As the price of this generality, however, we must observe that, when we are unwilling to disregard torsion phenomena, i.e. to tensor with $\mathbb Q$, the properties of $\Gr^\bullet(X)$ are less satisfactory in various respects than those of its two competitors (Exp. XIV no.~4). Thus, supposing again that $f:X\to Y$ is a proper morphism of noetherian regular schemes equipped with ample invertible Modules, it is not true in general that the homomorphism $f_*:K(X)\longrightarrow K(Y)$ maps $\operatorname{Filt}^i K(X)$ into $\operatorname{Filt}^{i-d} K(Y)$, where $d$ is the relative dimension of $X$ over $Y$, and hence defines, by passage to the associated graded rings, a homomorphism of degree $-d$ from $\Gr^\bullet(X)$ to $\Gr^\bullet(Y)$. However, we can prove (Exp. VII) that this holds after tensoring with $\mathbb Q$, so that we obtain a homomorphism of degree $-d$
\sgaStableTarget{sga6.e00.eq.2-5}\begin{equation}\tag{2.5}
f_*:\Gr^\bullet(X)_{\mathbb Q}\longrightarrow \Gr^\bullet(Y)_{\mathbb Q},
\end{equation}
playing the role of the homomorphism ``direct image of cycles'' $A(X)\to A(Y)$.

With this understood, we see that we have all the structural elements needed to give a meaning to formula (1.2) without postulating the existence of a base field, $f:X\to Y$ being a proper morphism of noetherian regular schemes admitting ample invertible Modules.

